\hspace{3mm}

\centerline{\bf \Huge Тренировочная Олимпиада}

\hspace{3mm}

\problem{1} В озере водятся караси, окуни и щуки. Два рыбака поймали вместе 70 рыб, причём $\dfrac{5}{9}$ улова первого рыбака составляли караси, а $\dfrac{7}{17}$ улова второго — окуни. При этом первый поймал столько же карасей, сколько второй, и столько же окуней, сколько второй. Сколько щук поймал первый рыбак и сколько — второй.

\hspace{3mm}

\problem{2} По кругу стоят 22 человека, Каждый из них — рыцарь (кото- рый всегда говорит только правду) или лжец (который всегда лжет). Каждый из них произнес фразу: «Следующие 10 человек по часовой стрелке после меня — лжецы». Сколько среди этих 22 людей лжецов?

\hspace{3mm}

\problem{3} Можно ли пронумеровать грани куба числами 1, 2, 3, 4, 5 и 6 так, чтобы номер каждой грани был делителем суммы номеров соседних граней? Если да — как, если нет — почему?

\hspace{3mm}

\problem{4} Когда Винни-Пух пришел в гости к Кролику, он съел 3 тарелки мёда, 4 тарелки сгущёнки и 2 тарелки варенья, а после этого не смог выйти наружу из-за того, что сильно растолстел от такой еды. Но известно, что если бы он съел 2 тарелки мёда, 3 тарелки сгущёнки и 4 тарелки варенья или 4 тарелки мёда, 2 тарелки сгущёнки и 3 тарелки варенья, то спокойно смог бы покинуть нору гостеприимного Кролика. От чего больше толстеют: от варенья или от сгущёнки?

\hspace{3mm}

\problem{5} В каждой клетке клетчатой доски размером 50×50 записано по числу. Известно, что каждое число в 3 раза меньше суммы всех чисел, записанных в клетках, соседних с ним по стороне, и в 2 раза меньше суммы всех чисел, записанных в клетках, соседних с ним по диагонали. Докажите, что каждую клетку доски можно покрасить в красный или синий цвет так, что сумма всех чисел, записанных в красных клетках, равна сумме всех чисел, записанных в синих клетках.